%% A projected reflected gradient method with feasible inexact projections:
%% a computable error certificate and a cost analysis.
%%
%% Target: Optimization (Q2) or Journal of Applied and Numerical Optimization.
%% Compile: pdflatex main && bibtex main && pdflatex main && pdflatex main
\documentclass[11pt]{article}

\usepackage[margin=1in]{geometry}
\usepackage{amsmath,amssymb,amsthm}
\usepackage{booktabs}
\usepackage{graphicx}
\usepackage[hidelinks]{hyperref}

\theoremstyle{plain}
\newtheorem{theorem}{Theorem}[section]
\newtheorem{lemma}[theorem]{Lemma}
\newtheorem{proposition}[theorem]{Proposition}
\newtheorem{corollary}[theorem]{Corollary}
\theoremstyle{definition}
\newtheorem{definition}[theorem]{Definition}
\newtheorem{assumption}{Assumption}
\theoremstyle{remark}
\newtheorem{remark}[theorem]{Remark}

\newcommand{\R}{\mathbb{R}}
\newcommand{\HH}{\mathcal{H}}
\newcommand{\inner}[2]{\langle #1, #2 \rangle}
\newcommand{\norm}[1]{\lVert #1 \rVert}
\newcommand{\proj}[1]{P_{#1}}
\newcommand{\TV}{\mathrm{TV}}

\title{A projected reflected gradient method with feasible inexact
projections: a computable error certificate and a cost analysis}

\author{Dao Trong Hieu\thanks{Corresponding author. E-mail:
tronghieu.ghe@gmail.com}}

\date{\today}

\begin{document}
\maketitle

\begin{abstract}
We study the projected reflected gradient (PRG) method of Malitsky for
monotone variational inequalities in the practically relevant regime where the
projection onto the feasible set has no closed form and must be computed by an
inner solver, so that only an \emph{inexact} projection is available. Our
contribution is threefold. First, we analyse PRG with feasible inexact
projections and show that weak convergence is preserved under a summable error
sequence; the argument combines the one-step inequality of PRG with the
bounded-perturbation-resilience mechanism, and we state explicitly that the
convergence statement is an extension of known results rather than a new
mechanism. Second, and this is the practically useful part, we give a
\emph{computable} error certificate for the inexact projection: when the
projection subproblem is strongly convex, the duality gap of the subproblem
bounds the distance to the exact projection, which yields a stopping rule that
an implementation can actually check. Existing inexact-projection schemes
either assume the error sequence abstractly or rely on knowledge of the exact
projection, which is precisely the quantity the method is designed to avoid.
Third, we use the certificate to define an adaptive inner budget and we give a
cost analysis together with an extensive numerical study on projections onto a
total-variation ball, where the inner solver is a Chambolle--Pock iteration
with dual warm starting. On image deblurring problems the adaptive scheme
reaches a prescribed variational residual level with $13$--$19$ times fewer
inner iterations and $5$--$10$ times less algorithm wall-clock time than the
certified exact projection, while returning strictly feasible iterates. We
report both cost measures, we fix the target residual levels in advance, and we
tune the baseline with the same budget as the proposed scheme; we also document
the configurations in which the adaptive scheme \emph{loses}.
\end{abstract}

\noindent\textbf{Keywords:} variational inequality; projected reflected
gradient; inexact projection; duality gap; total variation; image deblurring.

\noindent\textbf{MSC (2020):} 47J20; 65K15; 90C25; 68U10.

\section{Introduction}
\label{sec:intro}

Let $\HH$ be a real Hilbert space, $D \subseteq \HH$ a nonempty closed convex
set and $F : \HH \to \HH$ an operator. The variational inequality problem
$\mathrm{VI}(F,D)$ is
\begin{equation}
\label{eq:vi}
\text{find } x^\ast \in D \quad\text{such that}\quad
\inner{F(x^\ast)}{x - x^\ast} \ge 0 \quad \text{for all } x \in D .
\end{equation}
We write $S$ for the solution set and assume throughout that $S \neq
\emptyset$.

Projection-type methods for \eqref{eq:vi} are classical. Among them the
projected reflected gradient (PRG) method of Malitsky \cite{Malitsky2015} is
attractive because it needs only \emph{one} projection and \emph{one} operator
evaluation per iteration, in contrast with extragradient-type schemes that need
two. Its iteration is
\begin{equation}
\label{eq:prg}
x^{k+1} = \proj{D}\bigl( x^k - \lambda F(2x^k - x^{k-1}) \bigr),
\end{equation}
and it converges weakly for monotone Lipschitz $F$ provided $\lambda \in
(0,(\sqrt{2}-1)/L)$.

\paragraph{The gap this paper addresses.}
Analyses of \eqref{eq:prg} assume that $\proj{D}$ is available exactly. In the
applications that motivate constrained variational inequalities this is false.
A representative case, and the one we study numerically, is the total-variation
ball $D = \{x : \TV(x) \le \tau\}$ arising in image reconstruction: projecting
onto $D$ has no closed form and requires an inner iterative solver, so every
outer iteration works with an inexact projection. Two questions follow, and the
literature answers neither for PRG:
\begin{enumerate}
\item Does PRG still converge when $\proj{D}$ is replaced by an inexact
      projection with controlled error?
\item How does one \emph{check}, inside an implementation, that the error
      tolerance has been met, without knowing the exact projection?
\end{enumerate}
Question~2 is the more consequential one and is usually left open. Inexact
projection schemes such as \cite{DiazMillan2024} control the error through a
relative criterion, and abstract analyses assume a summable error sequence
$\sum_k \varepsilon_k < \infty$; in both cases the practical stopping rule for
the inner solver is either assumed or requires the exact projection, which is
the very object the outer method avoids computing.

\paragraph{Contributions and an honest positioning.}
\begin{itemize}
\item \emph{Convergence (Section~\ref{sec:conv}).} We show that PRG with
      feasible inexact projections and a summable error sequence retains weak
      convergence for monotone Lipschitz operators. We state plainly that this
      is an extension obtained by combining the one-step inequality of
      \cite{Malitsky2015} with the bounded-perturbation-resilience mechanism of
      \cite{Censor2017}, and not a new mechanism; the technical content lies in
      the fact that PRG is not Fej\'er monotone, so the perturbation must be
      carried through the reflection step, where an error at the base point is
      amplified by the factor two of the reflection and enters through the
      projection inequality of the previous step with a factor $1/\lambda$.
\item \emph{A computable certificate (Section~\ref{sec:cert}).} The projection
      subproblem $\min_x \tfrac12 \norm{x-v}^2$ subject to $x \in D$ is
      $1$-strongly convex, hence
      $\norm{x - \proj{D}(v)}^2 \le 2\,\mathrm{gap}(x,p)$, where the gap is
      computed from the primal--dual pair produced by the inner solver. This
      turns ``$\varepsilon_k$-inexact projection'' from an assumption into a
      checkable stopping rule. The inequality itself is standard convex
      analysis; the contribution is its use as the interface between the outer
      convergence theory and the inner solver, which is exactly what is missing
      in the inexact-projection literature.
\item \emph{Adaptive budget and cost analysis (Sections~\ref{sec:cost}
      and~\ref{sec:num}).} The certificate defines an adaptive inner budget:
      iterate the inner solver until the certificate meets a prescribed
      schedule $\varepsilon_k$. We analyse the resulting inner cost and report
      an extensive numerical study.
\end{itemize}

\paragraph{What this paper does not claim.}
We do not claim a new convergence mechanism, and we do not claim strong
convergence: adding an anchoring step to obtain strong convergence breaks the
projection chain that the one-step inequality of PRG relies on, and our numerical
evidence in Section~\ref{sec:num} shows that anchoring also pushes the
summability of the error sequence to the boundary of divergence. We also do not
claim an advantage in reconstruction quality over learned or plug-and-play
methods; the advantage we measure is in \emph{cost at a fixed accuracy}.

\section{Preliminaries}
\label{sec:prelim}

We use the following standing assumptions.

\begin{assumption}
\label{as:D}
$D \subseteq \HH$ is nonempty, closed and convex, and $S \neq \emptyset$.
\end{assumption}

\begin{assumption}
\label{as:F}
$F$ is monotone on $\HH$, i.e.\ $\inner{F(u)-F(v)}{u-v}\ge 0$ for all $u,v \in
\HH$, and Lipschitz continuous on $\HH$ with constant $L$.
\end{assumption}

Monotonicity is required on all of $\HH$ rather than on $D$ only, because the
reflected point $2x^k - x^{k-1}$ generally lies outside $D$; the same
requirement appears in \cite{Malitsky2015}. In our numerical setting
$F(x) = B^{\!\top}(Bx-y)$ is affine with $B^{\!\top}B \succeq 0$, hence
monotone and Lipschitz on $\HH$.

\begin{assumption}
\label{as:step}
$\lambda \in \bigl(0, (\sqrt{2}-1)/L\bigr)$.
\end{assumption}

We recall two standard facts. For $v \in \HH$ and $z \in D$,
\begin{equation}
\label{eq:projchar}
\inner{v - \proj{D}(v)}{z - \proj{D}(v)} \le 0 ,
\end{equation}
and for $a,b \in \HH$, $2\inner{a}{b} = \norm{a}^2 + \norm{b}^2 -
\norm{a-b}^2$.

\begin{lemma}[quasi-Fej\'er, \cite{Combettes2001}]
\label{lem:qf}
Let $\{s_k\}$ be nonnegative with $s_{k+1} \le s_k + \delta_k$ and $\sum_k
\delta_k < \infty$. Then $\{s_k\}$ converges.
\end{lemma}

\section{The algorithm}
\label{sec:alg}

\begin{definition}[feasible inexact projection]
\label{def:inexact}
Given $v \in \HH$ and $\varepsilon \ge 0$, a point $w$ is an
$\varepsilon$-inexact feasible projection of $v$ onto $D$, written $w \in
\proj{D}^{\varepsilon}(v)$, if $w \in D$ and $\norm{w - \proj{D}(v)} \le
\varepsilon$.
\end{definition}

Feasibility of $w$ is not cosmetic. It is what allows us to use
$\inner{F(z)}{x^k - z} \ge 0$ for $z \in S$ in Lemma~\ref{lem:tele} below;
without it the telescoping term loses its sign. Section~\ref{sec:cert}
shows how the inner solver produces feasible iterates and certifies
$\varepsilon$.

\begin{center}
\fbox{\begin{minipage}{0.92\textwidth}
\textbf{Algorithm 1 (PRG with feasible inexact projections).}\\[2pt]
Choose $x^{-1} = x^0 \in D$, $\lambda$ as in Assumption~\ref{as:step}, and a
sequence $\{\varepsilon_k\}$ with $\sum_k \varepsilon_k < \infty$.\\
For $k = 0,1,2,\dots$:
\begin{align*}
r^k     &= 2x^k - x^{k-1}, \\
x^{k+1} &\in \proj{D}^{\varepsilon_k}\bigl( x^k - \lambda F(r^k) \bigr).
\end{align*}
\end{minipage}}
\end{center}

We write $\bar x^{k+1} = \proj{D}(x^k - \lambda F(r^k))$ for the exact
projection, so that $\norm{x^{k+1} - \bar x^{k+1}} \le \varepsilon_k$, and
$e^k = x^{k+1} - \bar x^{k+1}$.

\section{Convergence}
\label{sec:conv}

\begin{lemma}[one-step inequality]
\label{lem:onestep}
For every $z \in D$,
\[
\norm{\bar x^{k+1} - z}^2 \le \norm{x^k - z}^2 - \norm{x^k - \bar x^{k+1}}^2
+ 2\lambda \inner{F(r^k)}{z - \bar x^{k+1}} .
\]
\end{lemma}

\begin{proof}
Apply \eqref{eq:projchar} with $v = x^k - \lambda F(r^k)$ and the projection
$\bar x^{k+1}$, which gives
$\inner{x^k - \bar x^{k+1}}{z - \bar x^{k+1}} \le \lambda
\inner{F(r^k)}{z - \bar x^{k+1}}$. Expand the left-hand side with the
polarisation identity, using $(x^k - \bar x^{k+1}) - (z - \bar x^{k+1}) = x^k
- z$, and rearrange.
\end{proof}

\begin{lemma}[monotonicity step]
\label{lem:mono}
For every $z \in S$,
\[
2\lambda \inner{F(r^k)}{z - \bar x^{k+1}} \le
2\lambda \inner{F(z)}{z - r^k} + 2\lambda \inner{F(r^k)}{r^k - \bar x^{k+1}} .
\]
\end{lemma}

\begin{proof}
Split $\inner{F(r^k)}{z - \bar x^{k+1}} = \inner{F(r^k)}{z-r^k} +
\inner{F(r^k)}{r^k - \bar x^{k+1}}$ and apply Assumption~\ref{as:F} to the
pair $(r^k, z)$, which yields $\inner{F(r^k)}{z-r^k} \le
\inner{F(z)}{z-r^k}$. This is the only place where monotonicity is used, and
it is used at a pair whose first entry lies outside $D$.
\end{proof}

\begin{lemma}[telescoping of the solution term]
\label{lem:tele}
Let $z \in S$ and set $a_k = \inner{F(z)}{z - x^k}$. If $x^k \in D$ for all
$k$, then $a_k \le 0$ and
\[
\inner{F(z)}{z - r^k} = 2a_k - a_{k-1} .
\]
\end{lemma}

\begin{proof}
$a_k \le 0$ is the definition of $z \in S$ together with $x^k \in D$. The
identity follows from $z - r^k = 2(z-x^k) - (z-x^{k-1})$.
\end{proof}

\begin{lemma}[uniform perturbation]
\label{lem:pert}
For every $z \in \HH$,
$\norm{x^{k+1} - z} \le \norm{\bar x^{k+1} - z} + \varepsilon_k$ and
$\norm{x^{k+1} - z}^2 \le \norm{\bar x^{k+1} - z}^2
 + 2\varepsilon_k \norm{\bar x^{k+1} - z} + \varepsilon_k^2$.
\end{lemma}

\begin{theorem}
\label{thm:weak}
Let Assumptions~\ref{as:D}--\ref{as:step} hold, let $\{x^k\}$ be generated by
Algorithm~1 with $\sum_k \varepsilon_k < \infty$, and assume $\HH$ is finite
dimensional or $F$ is sequentially weakly continuous on $D$. Then $\{x^k\}$
converges weakly to a point of $S$.
\end{theorem}

\begin{proof}[Proof sketch and status]
Combining Lemmas~\ref{lem:onestep}--\ref{lem:tele} and following the energy
argument of \cite[Lemma~9 and Theorem~4.2]{Malitsky2015} with the potential
\[
\varphi_k = \norm{x^k-z}^2 + \lambda L \norm{x^k - r^{k-1}}^2 - 2\lambda
a_{k-1},
\]
one obtains $\varphi_{k+1} \le \varphi_k + 2\lambda a_k - c\,\norm{x^k - \bar
x^{k+1}}^2 + \delta_k$ with $c > 0$ under Assumption~\ref{as:step}. The
perturbation of Lemma~\ref{lem:pert} contributes $\delta_k = O(\varepsilon_k)$
on bounded sequences, and the error at the \emph{base point}, $x^k = \bar x^k +
e^{k-1}$, contributes further cross terms which are $O(\varepsilon_{k-1}/
\lambda)$ because the projection inequality of the previous step is divided by
$\lambda$, and which are amplified by the factor two of the reflection. All of
these are summable, so Lemma~\ref{lem:qf} applies to $\varphi_k$; since $a_k
\le 0$, summing also gives $\sum_k (-a_k) < \infty$ and hence $a_k \to 0$,
from which $\norm{x^k - z}$ converges. The negative term gives $\norm{x^k -
\bar x^{k+1}} \to 0$ and hence $\norm{x^k - x^{k-1}} \to 0$; weak cluster
points then belong to $S$ by Minty's lemma, and Opial's lemma concludes.

\emph{Status.} The four lemmas above are complete. The assembly of the
potential $\varphi_k$ carrying the base-point error, and the two limit steps at
the end, are written here at the level of a proof sketch and are the part of
the argument that must be verified line by line before submission. We flag this
explicitly rather than presenting the theorem as fully established.
\end{proof}

\begin{remark}[honest positioning]
\label{rem:position}
Theorem~\ref{thm:weak} should be read as an extension, not as a new mechanism.
If $\{\varepsilon_k\}$ is a prescribed summable sequence, the statement is
close to what one obtains by combining \cite{Malitsky2015} with the
bounded-perturbation-resilience results of \cite{Censor2017}; the reason it is
not an immediate corollary is that PRG is not Fej\'er monotone and the
perturbation enters through the reflection, as described in the proof. The
practical value of this paper lies in Sections~\ref{sec:cert}--\ref{sec:num}.
\end{remark}

\section{A computable error certificate}
\label{sec:cert}

Definition~\ref{def:inexact} is not directly checkable: verifying $\norm{w -
\proj{D}(v)} \le \varepsilon$ presupposes $\proj{D}(v)$. We now remove this
circularity for the case where the projection subproblem is solved by a
primal--dual method.

Consider $D = \{x : \TV(x) \le \tau\} \cap [0,1]^n$ and the subproblem
\begin{equation}
\label{eq:sub}
\proj{D}(v) = \arg\min_x \; G(x) + F_{\mathrm{ind}}(Kx), \qquad
G(x) = \tfrac12\norm{x-v}^2 , \quad K = \nabla ,
\end{equation}
with $F_{\mathrm{ind}}$ the indicator of the $\ell_{2,1}$ ball of radius
$\tau$. Chambolle--Pock \cite{ChambollePock2011} produces a primal--dual pair
$(x,p)$. Since $G$ is $1$-strongly convex,
\begin{equation}
\label{eq:cert}
\norm{x - \proj{D}(v)}^2 \le 2\bigl( P(x) - P(\proj{D}(v)) \bigr)
\le 2\bigl( P(x) - Dl(p) \bigr) =: 2\,\mathrm{gap}(x,p) ,
\end{equation}
where $P(x) = \tfrac12\norm{x-v}^2$ for feasible $x$ and
$Dl(p) = \inner{v}{-\mathrm{div}\,p} - \tfrac12\norm{\mathrm{div}\,p}^2 -
\tau\norm{p}_{2,\infty}$.

\begin{proposition}
\label{prop:cert}
$\sqrt{2\,\mathrm{gap}(x,p)}$ is a computable upper bound for $\norm{x -
\proj{D}(v)}$, evaluated from the primal--dual pair alone.
\end{proposition}

Inequality \eqref{eq:cert} is standard convex analysis. Its role here is what
matters: it is the missing interface between the outer convergence theory,
which speaks of $\varepsilon_k$, and the inner solver, which can only measure
its own iterates. Combined with a feasibility step (rescaling towards the mean,
which preserves membership in $D$ since $\TV$ is positively homogeneous about
constants), it produces exactly the object required by
Definition~\ref{def:inexact}.

\begin{remark}[what does not work]
\label{rem:fail}
It is tempting to replace Definition~\ref{def:inexact} by the relaxed
projection inequality of \cite{DiazMillan2024}, in which the error enters
quadratically and a constant relative tolerance suffices. We report that this
is not implementable here. Checking it requires the support function of $D$;
bounding it by the enclosing box is far too loose (in our experiments the left
side stays at $1.2 \times 10^{2}$ while the right side is $3\times10^{-1}$,
even after $2000$ inner iterations), and deriving it from the certificate
\eqref{eq:cert} costs a factor: the error enters the relaxed inequality
linearly while the criterion is quadratic in the displacement, so one would
need $\varepsilon \sim \norm{w-u}^2$, i.e.\ five orders of magnitude below the
displacement at $k=200$ in our runs. We record this to save others the attempt.
\end{remark}

\section{Adaptive inner budget and cost}
\label{sec:cost}

The certificate turns the schedule $\{\varepsilon_k\}$ into an implementable
rule: at outer iteration $k$, run the inner solver, with dual warm starting
from iteration $k-1$, until $\sqrt{2\,\mathrm{gap}} \le \varepsilon_k$. We take
$\varepsilon_k = \varepsilon_0/(k+1)^p$ with $p>1$, which is summable as
required by Theorem~\ref{thm:weak}.

Two costs must be reported and we report both: the total number of inner
iterations, and the wall-clock time of the algorithm proper, measured with GPU
synchronisation and excluding all instrumentation. The first alone is
misleading, because an inner iteration of the adaptive scheme also evaluates
the certificate and is therefore more expensive than a bare Chambolle--Pock
step.

\begin{remark}[price of a checkable rule]
The certificate is conservative: its ratio to the true error grows from about
$2$ to about $12$ as the iterates approach the solution, and stopping by the
certificate costs $2.6$--$5.8$ times more inner iterations than stopping by the
true (unknowable) error. This is the price of a rule that an implementation can
actually check.
\end{remark}

\section{Numerical experiments}
\label{sec:num}

\paragraph{Setup.}
Image deblurring $y = Bx + \epsilon$ with $B$ a normalised blur (Gaussian
$9\times9$, $\sigma=1.6$; and motion, length $9$, $30^\circ$), noise level
$0.05$, grayscale patches of side $96$, eight test images, $K=150$ outer
iterations, $F(x)=B^{\!\top}(Bx-y)$, $D$ a TV ball with radius $\tau = 0.55\,
\TV(x_{\mathrm{true}})$, and $\lambda = 0.9(\sqrt2-1)/L$ with $L =
\norm{B^{\!\top}B}$ estimated by power iteration. Runs are on an NVIDIA L40S.

\paragraph{Protocol.}
Target residual levels are fixed \emph{in advance} and independently of any
configuration. Both groups use the same certificate, so neither side is given
information the other lacks. The baseline (certified exact projection) is tuned
over the same number of configurations (eight) as the proposed scheme.

\paragraph{Cost at fixed accuracy.}
Table~\ref{tab:cost} reports, for each prescribed residual level, the cheapest
configuration of each group.

\begin{table}[t]
\centering
\caption{Cost to reach a prescribed variational residual, Gaussian blur. Time
is algorithm wall-clock, instrumentation excluded.}
\label{tab:cost}
\begin{tabular}{lrrrrrr}
\toprule
& \multicolumn{2}{c}{inner iterations} & \multicolumn{2}{c}{time (s)} &
\multicolumn{2}{c}{speed-up} \\
\cmidrule(lr){2-3}\cmidrule(lr){4-5}\cmidrule(lr){6-7}
residual & exact & adaptive & exact & adaptive & iter. & time \\
\midrule
$3.0\times10^{-2}$ & 3796 & 198 & 2.48 & 0.52 & 19.2 & 4.7 \\
$2.0\times10^{-2}$ & 3858 & 275 & 5.13 & 0.70 & 14.0 & 7.4 \\
$1.5\times10^{-2}$ & 5738 & 357 & 8.61 & 0.88 & 16.1 & 9.7 \\
$1.2\times10^{-2}$ & 8509 & 476 & 9.73 & 1.12 & 17.9 & 8.7 \\
$1.0\times10^{-2}$ & 8566 & 649 & 11.25 & 1.43 & 13.2 & 7.9 \\
\bottomrule
\end{tabular}
\end{table}

For motion blur the corresponding ranges are $17.2$--$23.6$ in inner iterations
and $6.2$--$8.7$ in time. All configurations return strictly feasible iterates
(constraint ratio $1.0000$), whereas a fixed inner budget violates the
constraint by up to a few percent.

\paragraph{Where the scheme loses.}
The schedule must not be too tight. With $\varepsilon_0 = 0.5$ and $p=1.5$ the
inner cost explodes to $5.2\times10^{5}$ iterations, about $50$ times
\emph{worse} than the exact projection; every configuration with $p=1.5$ fails.
The rule of thumb from our grid is $p$ slightly above $1$ and $\varepsilon_0$
large, around $2$ to $4$. We report this because a reader who picks the
schedule badly will observe the opposite of our headline numbers.

\paragraph{Summability of the error sequence.}
Theorem~\ref{thm:weak} requires $\sum_k \varepsilon_k < \infty$. With the
schedule above this holds by construction, but it is instructive to measure the
displacement $\norm{\bar x^{k+1}-x^k}$, which controls how tight the schedule
must be. Over $K=3000$ outer iterations its log--log slope is $-2.34$
(Gaussian) and $-2.84$ (motion), and the cumulative sums saturate (ratio $1.006$
and $1.000$ between half-horizon and horizon). If an anchoring step is added,
the slope degrades to $-1.01$ on motion blur, i.e.\ to the boundary of
divergence; at $K=600$ it still reads $-1.48$ and looks safe. This is direct
evidence for the design choice of Section~\ref{sec:alg}: anchoring would not
only break the projection chain used in Lemma~\ref{lem:onestep} but also push
summability to the edge.

\paragraph{A pseudomonotone, non-monotone instance.}
To confirm that the analysis is not vacuously restricted to a class larger than
the examples, we also run $F(x) = e^{-\norm{x}^2} A(x-b)$ on the unit ball in
$\R^{50}$ with $A \succeq 0$ singular. A numerical certificate finds pairs with
$\inner{F(u)-F(v)}{u-v} = -5.6\times10^{-6} < 0$, so $F$ is not monotone, while
the scheme still drives the residual from $5.0\times10^{-2}$ to
$1.25\times10^{-4}$.

\paragraph{Reproducibility.}
Code, raw traces and the analysis scripts are available at
\url{https://github.com/cloudysman/PIE_em_Hoang}.

\section{Conclusion}
\label{sec:concl}

We analysed the projected reflected gradient method with feasible inexact
projections, and, more importantly for practice, we closed the gap between the
$\varepsilon_k$ of the theory and what an implementation can measure, by using
the duality gap of the strongly convex projection subproblem as a computable
certificate. The resulting adaptive inner budget reaches a prescribed accuracy
with roughly an order of magnitude fewer inner iterations and $5$--$10$ times
less time than a certified exact projection, while keeping iterates feasible.
We have been explicit about the limits: the convergence statement is an
extension rather than a new mechanism, the certificate is conservative by a
factor of $2$--$12$, the relaxed-inequality route is not implementable here,
and a badly chosen schedule is worse than doing nothing.

Two directions remain. First, completing the assembly of the potential in
Theorem~\ref{thm:weak}, including the base-point error amplified by the
reflection; this is where whatever technical novelty the analysis carries is
concentrated. Second, obtaining strong convergence without destroying either
the projection chain or the summability margin, which our data suggests is not
merely a technical matter.

\begin{thebibliography}{9}

\bibitem{Malitsky2015}
Yu.~Malitsky.
\newblock Projected reflected gradient methods for monotone variational
inequalities.
\newblock \emph{SIAM Journal on Optimization}, 25(1):502--520, 2015.

\bibitem{Censor2017}
Y.~Censor, A.~Gibali, S.~Reich.
\newblock Bounded perturbation resilience of extragradient-type methods and
their applications.
\newblock \emph{Journal of Inequalities and Applications}, 2017:232, 2017.

\bibitem{DiazMillan2024}
R.~D\'iaz Mill\'an, O.~P.~Ferreira, J.~Ugon.
\newblock Extragradient method with feasible inexact projection to variational
inequality problem.
\newblock \emph{Computational Optimization and Applications}, 89:459--484,
2024.

\bibitem{ChambollePock2011}
A.~Chambolle, T.~Pock.
\newblock A first-order primal-dual algorithm for convex problems with
applications to imaging.
\newblock \emph{Journal of Mathematical Imaging and Vision}, 40:120--145, 2011.

\bibitem{Combettes2001}
P.~L.~Combettes.
\newblock Quasi-Fej\'erian analysis of some optimization algorithms.
\newblock In \emph{Inherently Parallel Algorithms in Feasibility and
Optimization and their Applications}, 115--152, 2001.

\bibitem{Tseng2000}
P.~Tseng.
\newblock A modified forward-backward splitting method for maximal monotone
mappings.
\newblock \emph{SIAM Journal on Control and Optimization}, 38(2):431--446,
2000.

\end{thebibliography}

\end{document}
